% Generated draft: requires mathematical and source review.
\documentclass[11pt]{article}
\usepackage[margin=25mm]{geometry}
\usepackage{fontspec}
\setmainfont{Latin Modern Roman}
\usepackage{amsmath,amssymb,mathtools,unicode-math}
\setmathfont{Latin Modern Math}
\usepackage{enumitem,needspace,xurl,hyperref}
\setlength{\parindent}{0pt}
\setlength{\parskip}{5pt}
\setlength{\emergencystretch}{3em}
\newcommand{\N}{\mathbb{N}}
\newcommand{\Z}{\mathbb{Z}}
\newcommand{\Q}{\mathbb{Q}}
\newcommand{\R}{\mathbb{R}}
\newcommand{\C}{\mathbb{C}}
\newcommand{\F}{\mathbb{F}}
\newcommand{\A}{\mathbb{A}}
\newcommand{\PP}{\mathbb{P}}
\newcommand{\E}{\mathbb{E}}
\newcommand{\eps}{\varepsilon}
\newcommand{\dd}{\,\mathrm{d}}
\begin{document}
\section*{UnsolvedMath ID: 700978}
\textbf{Source problem code:} catalogue entry
\subsection{Definitions and mathematical statement}
A group \(H\) has \emph{finite special rank} if there exists an integer \(r\geq 0\) such that every finitely generated subgroup of \(H\) can be generated by at most \(r\) elements. A group \(G\) is \emph{locally graded} if every nontrivial finitely generated subgroup \(K\leq G\) has a proper subgroup of finite index, equivalently, a nontrivial finite quotient. A group is \emph{periodic} if every element has finite order.

The statement \(G=AB\), for subgroups \(A,B\leq G\), means the setwise equality
\[
G=\{ab:a\in A,\ b\in B\}.
\]
No normality assumption on either factor is imposed.

\textbf{Question (N. S. Chernikov, Kourovka Notebook 13.60).} For every locally graded group \(G\) and all subgroups \(A,B\leq G\) satisfying \(G=AB\), does finite special rank of both \(A\) and \(B\) imply finite special rank of \(G\)? Explicitly, if there exist integers \(r_A,r_B\geq 0\) such that every finitely generated subgroup of \(A\), respectively \(B\), admits a generating set of size at most \(r_A\), respectively \(r_B\), must there exist an integer \(r_G\geq 0\) such that every finitely generated subgroup of \(G\) admits a generating set of size at most \(r_G\)? The question asks for existence of such a bound for each \(G\); no particular bound in terms of \(r_A\) and \(r_B\) is specified.

\textbf{Supplied qualification.} The source reports an affirmative answer under the additional assumption that both \(A\) and \(B\) are periodic, and attributes this case to the 1990 paper of N. S. Chernikov cited below. This reported result has not been independently verified here. The supplied registry labels the general question as open; that label is source metadata rather than an independently verified assessment.
\subsection{Short English statement}
If a locally graded group is the setwise product of two subgroups, each having a uniform finite bound on the number of generators of its finitely generated subgroups, must the whole group have such a bound? The source reports an affirmative answer when both factors are periodic.
\subsection{Sources}
\begin{itemize}
\item Local UnsolvedMath registry: supplied complete source record, permanent ID \(700978\), exact source code \texttt{catalogue entry}, title \emph{Kourovka Notebook Problem 13.60}. The record attributes the problem to N. S. Chernikov, gives the proposed year as 1995, and labels its status as open. Its recorded update date is \texttt{2026-10-16T14:00:00Z}; this is registry metadata, not independent source verification.
\item \emph{The Kourovka Notebook (Unsolved Problems in Group Theory)}, original problem 13.60, 13th issue (1995). The supplied background also identifies the 21st issue edition (2026), arXiv:1401.0300v48, \url{https://arxiv.org/abs/1401.0300v48}. The supplied collection description names E. I. Khukhro and V. D. Mazurov as editors. These source details are supplied attribution and have not been independently checked here.
\item N. S. Chernikov, \emph{Ukrain. Math. J.}, \textbf{42}, no. 7 (1990), 855--861. The supplied statement cites this paper for the affirmative answer when both factors are periodic; no article title, DOI, or article URL is supplied.
\end{itemize}
\end{document}
